\Needspace{25\baselineskip}
\section{Gauduchon metrics and the cofactor}
\label{sec:gauduchon-surface-uniqueness}

Integration against $w^2$ also measures the obstruction to using the unweighted cofactor. For smooth potentials, the error is exactly the pairing of $w^2$ with $\ddc T$. On a surface the Gauduchon condition makes this error vanish for every pair; in dimension three an additional term depends on the potentials.

\begin{proposition}\label{prop:mixed-ddc-uniqueness}
Let $n\geq2$ and let $u,v$ be smooth $\omega$-plurisubharmonic functions with $A^n=B^n=f\,\mathrm{d}V$ for a smooth $f\geq0$.
Then
\begin{equation}\label{eq:mixed-uniqueness-defect}
 \int_X \mathrm{d}w\wedge\dc w\wedge T
       =\frac12\int_X w^2\ddc T.
\end{equation}
If $\ddc T=0$ and $\{f>0\}$ is dense in $X$, then $u-v$ is constant.
\end{proposition}
\begin{proof}
By \eqref{uni-eq:linearization}, $\ddc w\wedge T=0$. Integrating \eqref{uni-eq:w2} against $T$ and applying Stokes twice yields
\[
 \int_X\ddc(w^2)\wedge T=\int_Xw^2\ddc T.
\]
Substitution proves \eqref{eq:mixed-uniqueness-defect}.

For the uniqueness conclusion, apply \cref{uni-thm:cofactor-intro}(1) with $h=1$. At every point with $f>0$, the form $B$ is positive definite, so the term $B^{n-1}$ makes the cofactor strictly positive. Its positivity locus is therefore dense, as required.
\end{proof}

\Needspace{12\baselineskip}
\begin{corollary}\label{cor:gauduchon-surface-uniqueness}
Let $X$ be a connected compact complex surface and let $\omega$ be a Gauduchon metric. If $f\geq0$ is smooth and its zero set has empty interior, then two smooth $\omega$-plurisubharmonic solutions, both normalized to have supremum zero, of
\[
 (\omega+\ddc u)^2=cf\,\mathrm{d}V
\]
with the same $c>0$ coincide.
\end{corollary}
\begin{proof}
In dimension two, $T=A+B$ and
\[
 \ddc T=2\ddc\omega=0,
\]
which is exactly the Gauduchon condition on a surface. Smooth nonnegative $f$ with zero set of empty interior is positive on a dense set. The proposition gives equality up to a constant; normalization removes that constant.
\end{proof}

The next section removes the Gauduchon condition on a surface by integrating a signed quadratic difference. A conformal change of metric would change the equation and therefore cannot provide that extension.

In dimension three, direct expansion gives
\[
 \ddc T=3\ddc(\omega^2)+3\ddc\omega\wedge\ddc(u+v).
\]
The Gauduchon condition cancels the first term but need not cancel the second. For a concrete check, on $X=(\mathbb C/(\mathbb Z+\mathrm{i}\mathbb Z))^3$, with coordinates $z_1,z_2,z_3$ on the universal cover, write $e_j=\mathrm{i}\,\mathrm{d}z_j\wedge \mathrm{d}\bar z_j$, $x_j=\operatorname{Re}z_j$, and set
\[
 H=2+\cos(2\pi x_1),\qquad
 \omega=e_1+H e_2+H^{-1}e_3.
\]
The expansion
$\omega^2=2H e_1\wedge e_2+2H^{-1}e_1\wedge e_3+2e_2\wedge e_3$
has $\ddc\omega^2=0$: differentiating either variable coefficient adds a repeated $e_1$. Take $u=v=\delta\cos(2\pi x_2)$ with nonzero $\delta$ sufficiently small. Then $A=B>0$, but
\[
 \ddc T=6\ddc\omega\wedge\ddc u\ne0,
\]
because the $e_1\wedge e_3$ coefficient of $\ddc\omega$ and the $e_2$ coefficient of $\ddc u$ are both nonzero at $x_1=x_2=0$. Thus a Gauduchon equation background in dimension three can have a cofactor with $\ddc T\ne0$, even for equal strictly positive solution forms.
